\chapter{The reals through Dedekind cuts}
\label{chap:reals-dedekind}

本章从 ordered set 的语言开始，说明 \(\Q\) 缺少 completeness，然后用
\textbf{Dedekind cuts} 构造 \(\R\)。正文保留 English mathematical terms：
total order, ordered field, upper bound, supremum, infimum, Dedekind cut,
decimal expansion, irrational number, sequence, limit, and Cauchy sequence。

Dedekind 的想法是：不要先假设 real number 已经存在，而是把一个 real number
理解为它在 \(\Q\) 中切出的左半边。若一个切口 \(A\subset\Q\) 代表某个数 \(r\)，
则 \(A\) 应当包含所有小于 \(r\) 的 rationals，不包含所有大于或等于 \(r\) 的
rationals，并且左半边没有最后一个 rational。这个写法把 ``gap'' 变成一个 set，
从而可以用 set-theoretic language 处理 order, operation, and completeness。

\section{Order language on \texorpdfstring{\(\Z\)}{Z} and \texorpdfstring{\(\Q\)}{Q}}

\begin{definition}[Total order]
设 \((X,\leq)\) 是 partially ordered set。若任意 \(x,y\in X\) 都满足
\[
  x\leq y \quad\text{or}\quad y\leq x,
\]
则称 \(\leq\) 是 total order 或 linear order。这里的重点是 comparability：
任意两个元素都可以放在同一条顺序线上比较。
\end{definition}

读 total order 时要分开两个层次。第一层是 partial order 已经提供的
reflexivity, antisymmetry, and transitivity；第二层是额外加入的 comparability。
许多集合本来有 inclusion order 或 divisibility order，却不是 total order，因为
存在两个元素无法比较。整数和有理数的 standard order 特别适合做分析，是因为每次
写出 \(x<y\)、\(x=y\)、\(x>y\) 时，三者中恰好有一种情况发生。

在 \(\Z\) 与 \(\Q\) 中，standard order 定义为
\[
  x\leq y \Longleftrightarrow 0\leq y-x.
\]
对 integers，这依赖每个 integer 恰好是 positive, zero, or negative。对 rationals，
同样用分子分母的符号定义 positivity。这个 order 与 addition 和 multiplication
相容，正是 later ordered field 的基础。

\begin{exerciseblock}[Lecture Notes Exercise 39]
Suppose \((X,\leq)\) is totally ordered. Then any \(Y\subset X\) equipped with
the restricted order is also totally ordered.
\end{exerciseblock}

\begin{solution}
Restricted order 定义为：对 \(y,y'\in Y\)，
\[
  y\leq_Y y' \Longleftrightarrow y\leq_X y'.
\]
因为 \(Y\subset X\)，任意 \(y,y'\in Y\) 也都是 \(X\) 的元素。由 \(X\) 上的
total order，可得
\[
  y\leq_X y' \quad\text{or}\quad y'\leq_X y.
\]
按 restricted order 的定义，这等价于
\[
  y\leq_Y y' \quad\text{or}\quad y'\leq_Y y.
\]
因此 \(Y\) 上的 restricted order 满足 total order 所需的 comparability。
Reflexivity, antisymmetry, and transitivity 已经从 partial order 的 restriction
继承，所以 \((Y,\leq_Y)\) 是 totally ordered set。
\end{solution}

\begin{explanation}
这个题的关键是不要额外证明 \(Y\) 有任何 algebraic closure。Total order 只讨论
元素之间是否可比较；当 \(Y\) 是 subset 时，两个 \(Y\) 中的元素仍然可以借用
\(X\) 中已有的比较结果。同类题的写法通常是：先写 restricted relation 的定义，
再把要证明的 order axiom 翻译回 ambient set \(X\) 中已经成立的 axiom。
\end{explanation}

\begin{exerciseblock}[Lecture Notes Exercise 40]
Prove the order compatibility properties in \(\Q\): if \(x\leq y\), then
\(x+z\leq y+z\); if \(x\geq0\) and \(y\geq0\), then \(xy\geq0\).
\end{exerciseblock}

\begin{solution}
由 standard order 的定义，\(x\leq y\) 意味着 \(y-x\geq0\)。于是
\[
  (y+z)-(x+z)=y-x\geq0,
\]
所以 \(x+z\leq y+z\)。

再设 \(x\geq0\) 且 \(y\geq0\)。若其中一个为 \(0\)，则 \(xy=0\geq0\)。若二者都
为 positive rationals，可写成
\[
  x=\frac{a}{b},\qquad y=\frac{c}{d},
\]
其中 \(a,b,c,d\in\Z\)，并且可取 \(a,b,c,d>0\)。因此
\[
  xy=\frac{ac}{bd},
\]
分子 \(ac\) 与分母 \(bd\) 都是 positive integers，所以 \(xy>0\)。合并零的情况，
可得 \(xy\geq0\)。
\end{solution}

\begin{explanation}
第一个性质说明 order 对 translation 不变；第二个性质说明 positive cone 在
multiplication 下封闭。这两个条件正是 ordered field definition 中 order 与
field operations 的连接点。处理 rational inequalities 时，先把 \(x\leq y\)
改写成 \(y-x\geq0\)，再用 positive rationals 的分子分母符号；这是把 order
problem 转成 arithmetic problem 的标准方法。
\end{explanation}

\begin{definition}[Ordered field]
一个 field \(F\) 若配有 total order \(\leq\)，并且满足
\[
  x\leq y\Rightarrow x+z\leq y+z,
  \qquad
  x\geq0,\ y\geq0\Rightarrow xy\geq0,
\]
则称为 ordered field。Lecture Notes Chapter 4 的目标是构造一个 complete
ordered field，也就是 \(\R\)。
\end{definition}

Ordered field 的定义把 algebra 和 order 绑在一起。第一条 compatibility 说：
平移不改变大小关系，所以 inequality 的两边可以同时加同一个数。第二条 compatibility
说：positive cone 对 multiplication 封闭，因此可以安全地把 non-negative
quantities 相乘。后面处理 \(\Q\)、Dedekind cuts、sequences 时，凡是要把
inequality 乘上一个正数、或从 \(x<y\) 推出 \(x+z<y+z\)，本质上都在使用这两条
ordered-field axioms。

\section{Bounds, supremum, and infimum}

设 \(Y\subset X\)，其中 \(X\) 是 partially ordered set。
\begin{itemize}
  \item \(m\in Y\) 是 maximum，若对所有 \(y\in Y\)，\(y\leq m\)。
  \item \(m\in Y\) 是 minimum，若对所有 \(y\in Y\)，\(m\leq y\)。
  \item \(u\in X\) 是 upper bound，若对所有 \(y\in Y\)，\(y\leq u\)。
  \item \(\ell\in X\) 是 lower bound，若对所有 \(y\in Y\)，\(\ell\leq y\)。
  \item \(s=\sup(Y)\) 是 supremum，若 \(s\) 是 upper bound，且任何 upper bound
  \(u\) 都满足 \(s\leq u\)。
  \item \(i=\inf(Y)\) 是 infimum，若 \(i\) 是 lower bound，且任何 lower bound
  \(\ell\) 都满足 \(\ell\leq i\)。
\end{itemize}

Maximum/minimum 和 supremum/infimum 的差别必须保持清楚。Maximum 必须属于
\(Y\)，而 supremum 只需要属于 ambient ordered set \(X\)。例如
\(\Q^+=\{q\in\Q:q>0\}\) 没有 minimum，因为 \(q/2\) 总是更小的 positive rational；
但它在 \(\Q\) 中有 infimum，数值是 \(0\)，即使 \(0\notin\Q^+\)。做题时可以按
两步检查 supremum：先证明候选 \(s\) 是 upper bound；再证明任何更小的数都不是
upper bound，或者等价地证明任何 upper bound \(u\) 都满足 \(s\leq u\)。

Supremum 和 infimum 若存在则 unique。若 \(s,s'\) 都是 \(Y\) 的 supremum，则
\(s\) 是 upper bound，而 \(s'\) 是 least upper bound，所以 \(s'\leq s\)。反过来
也有 \(s\leq s'\)，由 antisymmetry 得 \(s=s'\)。Infimum 的 proof 把 inequality
方向反转即可。

\begin{exerciseblock}[Lecture Notes Exercise 41]
Show that \(\Z\) has no maximum or minimum.
\end{exerciseblock}

\begin{solution}
先证明没有 maximum。假设 \(M\in\Z\) 是 maximum。因为 \(M+1\in\Z\)，而
\[
  M<M+1,
\]
这与 maximum 的定义矛盾：maximum 必须大于或等于每个 integer。因此 \(\Z\) 没有
maximum。

再证明没有 minimum。假设 \(m\in\Z\) 是 minimum。因为 \(m-1\in\Z\)，而
\[
  m-1<m,
\]
这与 minimum 的定义矛盾：minimum 必须小于或等于每个 integer。因此 \(\Z\) 没有
minimum。
\end{solution}

\begin{explanation}
Maximum 和 minimum 必须属于集合本身。Integers 可以一直加 \(1\) 或减 \(1\)，所以
任何候选元素都会被同一个集合里的另一个元素超过或低过。遇到类似的
``no maximum'' 或 ``no minimum'' 题时，proof strategy 是从一个任意候选元素开始，
在同一个 set 内构造一个严格更大或严格更小的元素。这样可以直接否定 universal
condition：maximum 要求 \(y\leq M\) 对所有 \(y\) 成立，而 \(M+1\in\Z\) 给出
反例；minimum 要求 \(m\leq y\) 对所有 \(y\) 成立，而 \(m-1\in\Z\) 给出反例。
\end{explanation}

\begin{proposition}[Infimum of positive rationals]
令 \(\Q^+=\{q\in\Q:q>0\}\)。则
\[
  \inf(\Q^+)=0.
\]
\end{proposition}

\begin{proof}
首先，\(0\) 是 lower bound，因为每个 \(q\in\Q^+\) 都满足 \(0<q\)。其次，若
\(\ell>0\)，则 \(\ell/2\in\Q^+\)，并且 \(\ell/2<\ell\)，所以 \(\ell\) 不是
lower bound。若 \(\ell\leq0\)，则 \(\ell\leq0\)。因此所有 lower bounds 都小于
或等于 \(0\)，而 \(0\) 本身是 lower bound，所以 \(0\) 是 greatest lower bound。
\end{proof}

\begin{definition}[Completeness]
一个 partially ordered set \(X\) 称为 complete，若每个 non-empty 且 bounded
above 的 subset 都有 supremum，并且每个 non-empty 且 bounded below 的 subset
都有 infimum。
\end{definition}

Completeness 是本章的主要分水岭。Order 本身只告诉我们怎样比较两个已经存在的
元素；completeness 额外保证某些由集合提出的 boundary 也在系统内部存在。对
calculus 而言，极限、nested intervals、least-upper-bound arguments 都需要这种
``boundary 不丢失'' 的性质。\(\Q\) 作为 ordered field 已经有良好的四则运算和
order compatibility，但仍会在平方根这样的 boundary 上漏掉元素。

\section{Why \texorpdfstring{\(\Q\)}{Q} is not complete}

考虑
\[
  S=\{x\in\Q:x^2<2\}.
\]
这个集合有 upper bounds，例如任何大于 \(2\) 的 rational 都是 upper bound；但它在
\(\Q\) 中没有 supremum。证明分三步：平方大于 \(2\) 的正候选还可以下降，平方小于
\(2\) 的候选还可以上升，而 \(\Q\) 中没有平方等于 \(2\) 的元素。

\(\Q\) 中不存在 \(s^2=2\) 的 rational，可按 parity 证明。若 \(s=a/b\) 是 lowest
terms 且 \(s^2=2\)，则 \(a^2=2b^2\)，所以 \(a^2\) 是 even，进而 \(a\) 是 even。
写 \(a=2k\)，代回得 \(4k^2=2b^2\)，即 \(b^2=2k^2\)，所以 \(b\) 也是 even。这与
\(a/b\) 已经约到 lowest terms 矛盾。因此 rational 候选 supremum 只能落在
\(s^2<2\) 或 \(s^2>2\) 两类，而这两类都会被下面的 perturbation arguments 排除。

\begin{exerciseblock}[Lecture Notes Exercise 42]
Show that a positive \(s\in\Q\) with \(s^2>2\) is an upper bound for
\(S=\{x\in\Q:x^2<2\}\).
\end{exerciseblock}

\begin{solution}
取任意 \(x\in S\)。若 \(x<0\)，由于 \(s>0\)，有 \(x<s\)。若 \(x\geq0\)，假设
\(x\geq s\)。两个数都 non-negative 时，order 与 multiplication 相容，所以
\[
  x^2\geq s^2>2.
\]
这与 \(x\in S\) 即 \(x^2<2\) 矛盾。因此在 \(x\geq0\) 的情况下也必须有 \(x<s\)。
任意 \(x\in S\) 都满足 \(x\leq s\)，所以 \(s\) 是 \(S\) 的 upper bound。
\end{solution}

\begin{explanation}
源文件中的排版/OCR 对 Exercise 42 的条件有歧义。若只写 \(s^2>2\) 而不要求
\(s>0\)，负数 \(s\) 不能成为 \(S\) 的 upper bound，因为 \(0\in S\) 且 \(0>s\)。
本章采用完成证明所需的 positive-upper-bound 版本。证明 strategy 是先分
\(x<0\) 与 \(x\geq0\)：负数自动小于 positive \(s\)，非负数则可以从
\(x\geq s\) 推到 \(x^2\geq s^2\)，再与 \(x^2<2<s^2\) 冲突。
\end{explanation}

\begin{exerciseblock}[Lecture Notes Exercise 43]
Suppose \(s^2<2\). Then there exists \(r\in\Q\), \(s<r\), such that
\(r^2<2\). Hence \(s\) is not an upper bound for \(S\).
\end{exerciseblock}

\begin{solution}
若 \(s<0\)，取 \(r=0\)。则 \(s<0=r\)，且 \(r^2=0<2\)。下面设 \(s\geq0\)。
令
\[
  M=2-s^2>0.
\]
由 positive rationals 的密度，可取 rational \(\eps>0\) 满足
\[
  \eps<1,\qquad \eps<\frac{M}{2s+1}.
\]
令 \(r=s+\eps\)。则 \(r\in\Q\) 且 \(r>s\)。并且
\[
  r^2=(s+\eps)^2=s^2+2s\eps+\eps^2
      <s^2+(2s+1)\eps
      <s^2+M=2.
\]
因此 \(r\in S\) 且 \(r>s\)，所以 \(s\) 不是 \(S\) 的 upper bound。
\end{solution}

\begin{explanation}
这个 proof 的难点是控制 \((s+\eps)^2\)。把剩余空间写成 \(M=2-s^2\)，再选择
\(\eps\) 使增加量 \(2s\eps+\eps^2\) 小于 \(M\)，就能把 \(s\) 往右推而仍留在
\(S\) 中。同类 perturbation 题都按这个模板写：先给 gap 命名，再展开新表达式，
最后把所有新增项压到 gap 以内。条件 \(\eps<1\) 的作用是把 \(\eps^2\) 控制在
\(\eps\) 级别，避免二次项单独失控。
\end{explanation}

\begin{exerciseblock}[Lecture Notes Exercise 44]
Suppose \(s>0\) and \(s^2>2\). Then there exists \(r\in\Q\), \(r<s\), which
is an upper bound for \(S\). Hence \(s\) is not a supremum of \(S\).
\end{exerciseblock}

\begin{solution}
令
\[
  M=s^2-2>0.
\]
取 rational \(\eps>0\) 满足
\[
  \eps<s,\qquad \eps<\frac{M}{2s}.
\]
令 \(r=s-\eps\)。则 \(0<r<s\)。同时
\[
  r^2=(s-\eps)^2=s^2-2s\eps+\eps^2
      >s^2-2s\eps
      >s^2-M=2.
\]
由 Lecture Notes Exercise 42，正 rational \(r\) 且 \(r^2>2\) 蕴含 \(r\) 是
\(S\) 的 upper bound。因为 \(r<s\)，\(s\) 不是 least upper bound。

若 \(S\) 在 \(\Q\) 中有 supremum \(t\)，则 \(t<0\) 不可能，因为 \(0\in S\)。
若 \(t^2<2\)，Exercise 43 排除它是 upper bound；若 \(t>0\) 且 \(t^2>2\)，本题
排除它是 least upper bound；而 \(t^2=2\) 在 \(\Q\) 中无解。因此 \(S\) 在 \(\Q\)
中没有 supremum。
\end{solution}

\begin{explanation}
Least upper bound 不能只是 upper bound，还要是所有 upper bounds 中最小的一个。
当 \(s^2>2\) 时，平方值距离 \(2\) 还有 \(M\) 的余量；选小的 \(\eps\) 可以把
\(s\) 左移到 \(r\)，同时保持 \(r^2>2\)，于是得到更小的 upper bound。这个题与
Exercise 43 成对出现：一个排除 ``太小的 boundary''，一个排除 ``太大的 boundary''，
二者合起来说明 missing supremum 正好卡在 irrational boundary。
\end{explanation}

\section{Dedekind cuts}

\begin{definition}[Dedekind cut: pair version]
一个 Dedekind cut 是一对 non-empty subsets \((A,B)\)，其中 \(A,B\subset\Q\)，并且
\[
  \Q=A\cup B,\qquad
  a<b\text{ for all }a\in A,\ b\in B,\qquad
  A\text{ has no maximum}.
\]
这里 \(A\) 是左半边，\(B\) 是右半边。
\end{definition}

Pair version 把 real number 看成 rational line 上的一刀。条件
\(a<b\text{ for all }a\in A,b\in B\) 表示左半边和右半边没有交错；\(\Q=A\cup B\)
表示没有 rational 被遗漏；\(A\) has no maximum 则排除把 boundary rational 放进
左半边的写法。若边界本来是 rational \(q\)，我们把 \(q\) 放在右半边，从而让左半边
始终是 open initial segment。

\begin{definition}[Dedekind cut: one-set version]
一个 subset \(A\subset\Q\) 称为 Dedekind cut，若：
\begin{enumerate}
  \item \(A\ne\emptyset\) 且 \(A\ne\Q\)；
  \item downward closed：若 \(x\in A\) 且 \(y<x\)，则 \(y\in A\)；
  \item no maximum：若 \(x\in A\)，则存在 \(y\in A\) 使 \(x<y\)。
\end{enumerate}
\end{definition}

One-set version 只保留左半边，因为右半边可以恢复为 \(\Q\setminus A\)。四个条件各有
作用：non-empty 让 cut 有内容，proper 让 cut 不是整个 rational line，downward
closed 让它真的是左半边，no maximum 让边界不被左半边吃掉。验证一个候选 cut 时，
不要把其中某一条当成其它条件的替代品；例如 downward closed 不会自动推出
proper，no maximum 也不会自动推出 non-empty。

\begin{exerciseblock}[Lecture Notes Exercise 45]
Prove that the pair version and one-set version of Dedekind cuts are equivalent.
\end{exerciseblock}

\begin{solution}
先设 \((A,B)\) 是 pair version 的 Dedekind cut。因为 \(A,B\) non-empty 且
\(\Q=A\cup B\)，所以 \(A\ne\emptyset\) 且 \(A\ne\Q\)。若 \(x\in A\) 且 \(y<x\)，
假设 \(y\notin A\)。由 \(\Q=A\cup B\)，有 \(y\in B\)。Pair version 要求每个
\(a\in A\) 都小于每个 \(b\in B\)，于是 \(x<y\)，这与 \(y<x\) 矛盾。所以
\(y\in A\)，即 \(A\) downward closed。Pair version 已经要求 \(A\) has no
maximum，因此 \(A\) 满足 one-set version。

反过来，设 \(A\subset\Q\) 满足 one-set version，令 \(B=\Q\setminus A\)。由
\(A\ne\emptyset\) 且 \(A\ne\Q\)，可得 \(A,B\) non-empty 且 \(\Q=A\cup B\)。取
\(a\in A\)、\(b\in B\)。若 \(b<a\)，由 downward closed 得 \(b\in A\)，与
\(b\in B\) 矛盾；若 \(b=a\)，同样与 \(A\cap B=\emptyset\) 矛盾。因此 \(a<b\)。
One-set version 的第三条给出 \(A\) has no maximum。所以 \((A,\Q\setminus A)\)
满足 pair version。
\end{solution}

\begin{explanation}
两个版本的差别只是是否显式写出右半边。One-set version 把右半边定义为 complement；
pair version 则直接写出 \(A\) 与 \(B\) 的 separation。Downward closed 正是由
``左边元素必须小于右边元素'' 推出来的。同类 equivalence proof 较可靠的写法是双向
证明：先从旧定义推出新定义的每一条，再从新定义重建旧定义中的对象。
\end{explanation}

\begin{definition}[Rational cut]
对 \(q\in\Q\)，定义
\[
  q^{\R}=\{x\in\Q:x<q\}.
\]
这个 cut 是 rational \(q\) 在 \(\R\) 中的 copy。若 pair cut \((A,B)\) 的 \(B\)
有 minimum，则该 cut 对应 rational number \(\min(B)\)。
\end{definition}

记号 \(q^{\R}\) 不是说 \(q\) 自己变成了一个 set，而是说 embedding
\(\Q\hookrightarrow\R\) 把 rational \(q\) 送到它的 rational predecessors。这个
embedding 保持 order：\(p<q\) 时 \(p^{\R}\subsetneq q^{\R}\)，因为所有小于 \(p\)
的 rationals 也小于 \(q\)，但中点 \((p+q)/2\) 只落在 \(q^{\R}\) 中。

\begin{exerciseblock}[Lecture Notes Exercise 46]
Show that for each \(q\in\Q\), the set \(q^{\R}\) is a Dedekind cut. Moreover,
show that \(i:\Q\to\R\), \(i(q)=q^{\R}\), is injective.
\end{exerciseblock}

\begin{solution}
令 \(q^{\R}=\{x\in\Q:x<q\}\)。它 non-empty，因为 \(q-1<q\)；它 proper，因为
\(q\notin q^{\R}\)。若 \(x\in q^{\R}\) 且 \(y<x\)，则 \(y<x<q\)，所以
\(y\in q^{\R}\)。若 \(x\in q^{\R}\)，则
\[
  x<\frac{x+q}{2}<q,
\]
并且 \((x+q)/2\in\Q\)，所以 \(q^{\R}\) has no maximum。故 \(q^{\R}\) 是
Dedekind cut。

若 \(i(p)=i(q)\)，即 \(p^{\R}=q^{\R}\)。假设 \(p<q\)。取
\[
  r=\frac{p+q}{2}.
\]
则 \(p<r<q\)，所以 \(r\in q^{\R}\) 但 \(r\notin p^{\R}\)，与两个 cuts 相等矛盾。
若 \(q<p\)，取 \(r=(p+q)/2\)，则 \(q<r<p\)，所以 \(r\in p^{\R}\) 但
\(r\notin q^{\R}\)，同样与 \(p^{\R}=q^{\R}\) 矛盾。因此 \(p=q\)，所以
\(i\) injective。
\end{solution}

\begin{explanation}
这里必须使用 strict inequality \(x<q\)，不能用 \(x\leq q\)。集合
\(\{x\in\Q:x\leq q\}\) 虽然 non-empty, proper, and downward closed，但有 maximum
\(q\)，会让同一个 rational boundary 出现不必要的二义性。证明 injective 时，若
两个 rational boundaries 不同，就在它们中间取一个 rational；这个中点只属于较大
boundary 的 cut，因而能分辨两个 cuts。
\end{explanation}

\section{Order and operations on cuts}

从现在起，\(\R\) 表示所有 one-set Dedekind cuts 的集合。Order 定义为
\[
  A\leq A' \Longleftrightarrow A\subseteq A'.
\]
这与 ``左半边越大，代表的 real number 越大'' 一致。Addition 定义为
\[
  A+A'=\{a+a':a\in A,\ a'\in A'\}.
\]
Negation 可以写成
\[
  -A=\{t\in\Q:\exists b\in\Q\setminus A\text{ such that }t<-b\}.
\]
这个公式在 rational 和 irrational 两种 cut 中都能处理边界；若 \(A=q^{\R}\)，
则 \(-A=(-q)^{\R}\)。

Positive cuts 的 multiplication 可以写成
\[
  A\cdot B=
  \{x\in\Q:x<0\text{ or }\exists a\in A,\exists b\in B
    \text{ with }a>0,\ b>0,\ x<ab\}.
\]
最后的 \(x<ab\) 是 downward closure 的写法：只要某个 positive product 在左边，
所有更小的 rationals 也必须在左边。一般 signs 的 multiplication 由 positive case
和 sign rules 定义，例如 \(A<0^{\R}\)、\(B>0^{\R}\) 时先计算 \((-A)B\)，再取
negative。用 cuts 做四则运算时，要同时检查两件事：结果仍是 Dedekind cut，并且
对 rational cuts 会回到熟悉的 rational arithmetic。

\begin{exerciseblock}[Lecture Notes Exercise 47]
If \(A=q^{\R}\) for some \(q\in\Q\), prove that \(-A=(-q)^{\R}\).
\end{exerciseblock}

\begin{solution}
使用 negation 的 cut formula：
\[
  -A=\{t\in\Q:\exists b\in\Q\setminus A\text{ such that }t<-b\}.
\]
当 \(A=q^{\R}\) 时，\(\Q\setminus A=\{b\in\Q:b\geq q\}\)。

若 \(t\in -A\)，则存在 \(b\geq q\) 使 \(t<-b\)。因为 \(-b\leq -q\)，所以
\[
  t<-b\leq -q,
\]
从而 \(t< -q\)，即 \(t\in(-q)^{\R}\)。

若 \(t\in(-q)^{\R}\)，则 \(t<-q\)。取 \(b=q\)，有 \(b\in\Q\setminus A\)，并且
\[
  t<-q=-b.
\]
所以 \(t\in -A\)。两边互相包含，故 \(-A=(-q)^{\R}\)。
\end{solution}

\begin{explanation}
Rational cut 的右半边有最小元素 \(q\)，所以求 negative 时边界变成 \(-q\)。公式中
用 \(t<-b\) 而不是 \(t\leq -b\)，正是为了保持 cut 的 no maximum 条件。做
negative 相关题目时，先找原 cut 的右半边，再把右半边取相反数；strict inequality
负责把新的 boundary 留在右半边。
\end{explanation}

\begin{exerciseblock}[Lecture Notes Exercise 48]
Prove that the Dedekind-cut construction of \(\R\) is complete.
\end{exerciseblock}

\begin{solution}
令 \(\mathcal S\subset\R\) non-empty，并假设它 bounded above。取一个 upper bound
\(U\in\R\)，于是每个 \(A\in\mathcal S\) 都满足 \(A\subseteq U\)。定义
\[
  A_*=\bigcup_{A\in\mathcal S}A.
\]
我们证明 \(A_*\) 是 Dedekind cut，并且是 \(\mathcal S\) 的 supremum。

因为 \(\mathcal S\) non-empty，取 \(A_0\in\mathcal S\)。由于 \(A_0\ne\emptyset\)，
有 \(A_*\ne\emptyset\)。又因为每个 \(A\subseteq U\)，所以 \(A_*\subseteq U\)。
而 \(U\ne\Q\)，因此 \(A_*\ne\Q\)。若 \(x\in A_*\) 且 \(y<x\)，则 \(x\in A\) 对某个
\(A\in\mathcal S\) 成立；由 \(A\) downward closed，得 \(y\in A\subseteq A_*\)。
若 \(x\in A_*\)，取 \(A\in\mathcal S\) 使 \(x\in A\)。因为 \(A\) has no maximum，
存在 \(z\in A\) 使 \(x<z\)，于是 \(z\in A_*\)。所以 \(A_*\) 是 Dedekind cut。

对每个 \(A\in\mathcal S\)，有 \(A\subseteq A_*\)，所以 \(A_*\) 是 upper bound。
若 \(V\) 是任意 upper bound，则每个 \(A\in\mathcal S\) 都满足 \(A\subseteq V\)，
所以
\[
  A_*=\bigcup_{A\in\mathcal S}A\subseteq V.
\]
因此 \(A_*\leq V\)，说明 \(A_*=\sup\mathcal S\)。

若 \(\mathcal S\) non-empty 且 bounded below，则集合
\[
  -\mathcal S=\{-A:A\in\mathcal S\}
\]
non-empty 且 bounded above。由上半部分，\(\sup(-\mathcal S)\) 存在。定义
\[
  \inf\mathcal S=-\sup(-\mathcal S).
\]
记 \(I=-\sup(-\mathcal S)\)。若 \(A\in\mathcal S\)，则 \(-A\leq\sup(-\mathcal S)\)，
反转 order 得 \(I\leq A\)，所以 \(I\) 是 lower bound。若 \(J\) 是任意 lower bound，
则 \(J\leq A\) 对所有 \(A\in\mathcal S\) 成立，反转 order 得
\(-A\leq -J\)，所以 \(-J\) 是 \(-\mathcal S\) 的 upper bound。由 supremum 的定义，
\(\sup(-\mathcal S)\leq -J\)，再反转 order 得 \(J\leq I\)。因此 \(I\) 是 greatest
lower bound。
\end{solution}

\begin{explanation}
Completeness 的核心不是对每个 cut 逐项计算，而是用 union。一个 bounded-above 的
cut family 的所有左半边合并后仍然是左半边；upper bound \(U\) 保证合并不会变成
整个 \(\Q\)。这就是 Dedekind construction 直接内建 supremum 的原因。遇到
Dedekind completeness 题时，supremum 通常先猜成 union；随后逐条验证
non-empty, proper, downward closed, and no maximum，再证明它是 least upper bound。
\end{explanation}

\begin{exerciseblock}[Worksheet 6 Exercise 1]
Determine which of the four listed subsets of \(\Q\) are Dedekind cuts.
\end{exerciseblock}

\begin{solution}
\begin{enumerate}[label=(\arabic*)]
  \item \(A=\{x\in\Q:x^3<17\}\) 是 Dedekind cut。它 non-empty，因为
  \(2^3=8<17\)；proper，因为 \(3^3=27>17\)。若 \(x\in A\) 且 \(y<x\)，则 cube
  function 在 \(\Q\) 上 strictly increasing，所以 \(y^3<x^3<17\)，故 \(y\in A\)。
  若 \(x^3<17\)，取 \(\eps>0\) rational 使
  \[
    \eps<1,\qquad \eps(3x^2+3|x|+1)<17-x^3.
  \]
  则 \((x+\eps)^3<17\)，所以 \(A\) has no maximum。

  \item \(B=\{x\in\Q:x^4<100\}\) 不是 Dedekind cut。因为 \(2\in B\)，但
  \(-4<2\)，而 \((-4)^4=256>100\)，所以 \(-4\notin B\)。这违反 downward closed。

  \item \(C=\{x\in\Q:x^2>-3\}\) 不是 Dedekind cut。每个 rational 都满足
  \(x^2\geq0>-3\)，所以 \(C=\Q\)。这违反 proper condition \(A\ne\Q\)。

  \item \(D=\{x\in\Q:x^5\leq2\}\) 是 Dedekind cut。它 non-empty，因为 \(1\in D\)；
  proper，因为 \(2^5=32>2\)。Fifth power 在 \(\Q\) 上 strictly increasing，所以
  downward closed 成立。若 \(x^5<2\)，可选小 rational \(\eps>0\) 使
  \((x+\eps)^5<2\)，从而得到更大的元素仍在 \(D\)。若 \(x^5=2\)，则把
  \(x=a/b\) 写成 lowest terms 会给出 \(a^5=2b^5\)，推出 \(2\mid a\)，继而
  \(2\mid b\)，与 lowest terms 矛盾。因此 \(D\) 没有 maximum。
\end{enumerate}
\end{solution}

\begin{explanation}
验证 Dedekind cut 时必须逐条检查 non-empty, proper, downward closed, and no
maximum。Polynomial inequality 的题目常见陷阱是偶次幂不保序；这正是
\(x^4<100\) 失败的原因。同类题先判断 defining function 是否 monotone：odd powers
在全体 rationals 上保序，even powers 通常要分正负区间处理。
\end{explanation}

\begin{exerciseblock}[Worksheet 6 Exercise 2]
For rational \(p,q\in\Q\), prove \(p^{\R}+q^{\R}=(p+q)^{\R}\).
\end{exerciseblock}

\begin{solution}
按 cut addition 的定义，
\[
  p^{\R}+q^{\R}
  =\{a+b:a<p,\ b<q,\ a,b\in\Q\}.
\]
若 \(x\in p^{\R}+q^{\R}\)，则 \(x=a+b\)，其中 \(a<p\) 且 \(b<q\)。相加得
\[
  x=a+b<p+q,
\]
所以 \(x\in(p+q)^{\R}\)。

反过来，若 \(x\in(p+q)^{\R}\)，即 \(x<p+q\)。令
\[
  \delta=\frac{p+q-x}{2}>0,\qquad a=p-\delta,\qquad b=x-a.
\]
则 \(a<p\)，而且
\[
  b=x-p+\delta
   =x-p+\frac{p+q-x}{2}
   =\frac{x-p+q}{2},
\]
所以
\[
  q-b
  =q-\frac{x-p+q}{2}
  =\frac{p+q-x}{2}
  =\delta>0.
\]
因此 \(b<q\)，且 \(x=a+b\)。于是 \(x\in p^{\R}+q^{\R}\)。两边互相包含，结论成立。
\end{solution}

\begin{explanation}
第二个 inclusion 的关键是把 \(x<p+q\) 留下的正 gap 平分成 \(\delta\)。这样可以
构造 \(a<p\) 和 \(b<q\)，并让它们的和正好等于 \(x\)。证明 cut operation 与
rational arithmetic 相容时，通常一个方向由 inequality 相加得到，另一个方向要把
目标 rational 拆成两个严格位于边界左侧的 pieces。
\end{explanation}

\begin{exerciseblock}[Worksheet 6 Exercise 3]
For a Dedekind cut \(A\), prove existence and uniqueness of \(-A\), prove the
rational case, and state the irrational formula.
\end{exerciseblock}

\begin{solution}
定义
\[
  N_A=\{t\in\Q:\exists b\in\Q\setminus A\text{ such that }t<-b\}.
\]
这个集合 non-empty：取 \(b\notin A\)，则 \(t=-b-1\in N_A\)。它 proper：取
\(a\in A\)。若 \(t\geq -a\)，则不可能存在 \(b\notin A\) 使 \(t<-b\)，因为这会给出
\(b<-t\leq a\)，再由 downward closed 得 \(b\in A\)，矛盾。若 \(t\in N_A\) 且
\(u<t\)，同一个 \(b\notin A\) 满足 \(u<t<-b\)，所以 \(u\in N_A\)。若 \(t<-b\)，
取 \(v=(t-b)/2\)。则 \(t<v<-b\)，所以 \(v\in N_A\)。故 \(N_A\) 是 Dedekind cut。

要说明 \(N_A+A=0^{\R}\)，先取 \(n\in N_A\)、\(a\in A\)。存在 \(b\notin A\) 使
\(n<-b\)。由于 \(a\in A\) 且 \(b\notin A\)，不能有 \(b\leq a\)，否则由 downward
closed 或 disjointness 得矛盾；因此 \(a<b\)。于是
\[
  n+a<-b+a<0,
\]
所以 \(N_A+A\subseteq0^{\R}\)。反向 inclusion 使用 cut 的 gap lemma：对任意
\(\eta>0\)，存在 \(a\in A\) 与 \(b\notin A\) 使 \(0<b-a<\eta\)。这个 lemma 可由
proper condition 和 no maximum 推出：若某个 \(\delta>0\) 满足
\(a+\delta\in A\) 对所有 \(a\in A\) 都成立，则从一个 \(a_0\in A\) 开始不断加
\(\delta\)，再用 Archimedean property 和 downward closure，会推出 \(A=\Q\)，与
proper condition 矛盾。因此存在 \(a\in A\) 使 \(a+\delta\notin A\)；取
\(\delta=\eta/2\)、\(b=a+\delta\) 即得 \(0<b-a<\eta\)。

给定 \(x<0\)，取 \(\eta=-x\)，再取这样的 \(a,b\)，令 \(n=x-a\)。因为
\[
  n=x-a<-b
  \Longleftrightarrow
  x<a-b,
\]
而 \(b-a<-x\) 等价于 \(x<a-b\)，所以 \(n\in N_A\)，且 \(x=n+a\)。因此
\(0^{\R}\subseteq N_A+A\)。
所以 \(N_A\) 是 additive inverse。

Uniqueness：若 \(B+A=0^{\R}\) 且 \(C+A=0^{\R}\)，则用 addition 的 associativity
和 zero law 得
\[
  B=B+0^{\R}=B+(A+C)=(B+A)+C=0^{\R}+C=C.
\]
所以 inverse unique。

若 \(A=q^{\R}\)，上一题已经给出 \(N_A=(-q)^{\R}\)。若 \(A\) 不是 rational cut，
则 \(\Q\setminus A\) 没有 minimum，此时
\[
  -A=\Q\setminus\{x\in\Q:-x\in A\}.
\]
\end{solution}

\begin{explanation}
求 negative 时最容易在 boundary 上出错。Rational cut 的右半边有最小元素，必须用
strict inequality 来避免 maximum；irrational cut 没有 rational boundary，所以
complement-reflection formula 与 general formula 合并。同类 inverse 题的检查顺序
是：先验证候选 inverse 自身是 cut，再证明 sum/product 的两个 inclusions，最后用
associativity 证明 uniqueness。
\end{explanation}

\begin{exerciseblock}[Worksheet 6 Exercise 4]
Let \(A=\{x\in\Q:x<0\text{ or }x^2<2\}\). Prove that \(A\) is not \(q^{\R}\)
for any \(q\in\Q\), and prove \(A\cdot A=2^{\R}\).
\end{exerciseblock}

\begin{solution}
先验证 \(A\) 是 Dedekind cut。它 non-empty，因为 \(1\in A\)；proper，因为
\(2\notin A\)。若 \(x\in A\) 且 \(y<x\)，当 \(x<0\) 时有 \(y<0\)，故 \(y\in A\)；
当 \(x\geq0\) 且 \(x^2<2\) 时，若 \(y<0\) 则 \(y\in A\)，若 \(0\leq y<x\) 则
\(y^2<x^2<2\)，也有 \(y\in A\)。No maximum 由 Exercise 43 的 perturbation 得到：
负数可向右移动一点仍为负；非负且平方小于 \(2\) 的元素可向右移动一点仍保持平方
小于 \(2\)。

若 \(A=q^{\R}\)，则 \(q\) 必须是 boundary。因为 \(1\in A\) 且 \(2\notin A\)，有
\(1<q\leq2\)。对所有 \(x<q\) 应有 \(x^2<2\)，对所有 \(x>q\) 应有 \(x^2>2\)。
这迫使 \(q^2=2\)。但 \(\Q\) 中不存在平方等于 \(2\) 的 rational，所以 \(A\) 不是
rational cut。

对 multiplication，\(A\) 是 positive cut。由 Dedekind multiplication 的 positive
definition，\(A\cdot A\) 包含所有 negative rationals；positive part 由 \(a,a'\in A\)
且 \(a,a'\geq0\) 的 products 产生。若 \(a,a'\geq0\) 且 \(a^2<2\)、\((a')^2<2\)，
则 \(aa'<2\)，所以 \(A\cdot A\subseteq2^{\R}\)。反过来，若 \(0\leq x<2\)，取
rational \(a\) 满足
\[
  0\leq x<a^2<2,
\]
这可由 Exercise 43 的上推 argument 应用于 \(\sqrt{x}\) 的 rational 近似得到。于是
\(a\in A\)，且 \(x<a\cdot a\)，所以 \(x\) 属于 \(A\cdot A\) 的 downward closure。
因此 \(A\cdot A=2^{\R}\)。
\end{solution}

\begin{explanation}
这个 cut 是 \(\sqrt2\) 的左半边。Irrationality 来自 boundary equation
\(q^2=2\) 在 \(\Q\) 中无解；乘法结论说明这个左半边在 Dedekind multiplication 下
确实平方成 \(2^{\R}\)。解决类似 ``polynomial cut'' 的题目时，先证明 cut 条件，
再证明它不是 rational cut，最后才处理 operation；这样 boundary 和 algebra 不会混在
同一步里。
\end{explanation}

\begin{exerciseblock}[Lecture Notes Exercise 50]
Show that there exists a positive real number \(r\in\R\) such that \(r^2=2\).
\end{exerciseblock}

\begin{solution}
令
\[
  r=\{x\in\Q:x<0\text{ or }x^2<2\}.
\]
Worksheet 6 Exercise 4 已经逐条证明 \(r\) 是 Dedekind cut，且 \(r\cdot r=2^{\R}\)。
它是 positive real number，因为 \(0^{\R}\subsetneq r\)：例如 \(1\in r\) 但
\(1\notin0^{\R}\)。因此存在 positive real \(r\) 满足 \(r^2=2^{\R}\)。按 usual
notation，称这个 \(r\) 为 \(\sqrt2\)。
\end{solution}

\begin{explanation}
这里的存在性不是从 decimal approximation 猜出来的，而是直接构造一个 cut。这个
cut 把所有 negative rationals 放进左边，再把 positive side 中平方小于 \(2\) 的
rationals 放进左边，从而形成正的 square root。
\end{explanation}

\begin{exerciseblock}[Homework 7 Problem 1]
Let \(A,B\in\R\). Determine which listed set operations necessarily produce
Dedekind cuts.
\end{exerciseblock}

\begin{solution}
\begin{enumerate}[label=(\alph*)]
  \item \(\{ab:a\in A,b\in B\}\) 不一定是 Dedekind cut。取
  \(A=B=0^{\R}=\{x\in\Q:x<0\}\)。则 \(1=(-1)(-1)\) 属于该集合，但
  \(1/2<1\) 不一定能写成两个 negative rationals 的 product；事实上 product set
  是 positive rationals，加上部分 large values，不包含所有更小的 rationals。
  Downward closed 失败。

  \item \(\{a-b:a\in A,b\in B\}\) 不一定是 Dedekind cut。仍取
  \(A=B=0^{\R}\)。对任意 \(q\in\Q\)，选 \(b<\min(0,-q)\)，令 \(a=q+b\)。则
  \(a<0\)、\(b<0\)，并且 \(q=a-b\)。所以该集合等于 \(\Q\)，违反 proper condition。

  \item \(\Q\setminus\{-x:x\in A\}\) 不一定是 Dedekind cut。取 \(A=q^{\R}\)。则
  \[
    \{-x:x\in A\}=\{y\in\Q:y>-q\},
  \]
  因而 complement 是 \(\{y\in\Q:y\leq -q\}\)。这个集合有 maximum \(-q\)，违反
  no maximum。

  \item \(A\cap B\) 一定是 Dedekind cut。Non-empty：取 \(a\in A\)、\(b\in B\)，
  再取 rational \(c<\min(a,b)\)，则 \(c\in A\cap B\)。Proper：若
  \(q\notin A\)，则 \(q\notin A\cap B\)。Downward closed 从 \(A\) 与 \(B\) 继承。
  若 \(x\in A\cap B\)，先在 \(A\) 中取 \(a_1>x\)，在 \(B\) 中取 \(b_1>x\)，再取
  rational \(y\) 满足
  \[
    x<y<\min(a_1,b_1).
  \]
  由 downward closed，\(y\in A\cap B\)。所以 \(A\cap B\) has no maximum。
\end{enumerate}
\end{solution}

\begin{explanation}
这个题训练的是 ``set operation'' 与 ``cut operation'' 的区别。Dedekind addition
和 multiplication 有专门定义；任意把 representatives 做 pointwise set operation
不会自动保持 cut 的四条性质。判断这类题时，先拿 rational cuts 或 \(0^{\R}\)
测试 closure；若要证明某个 operation 总能产生 cut，则必须回到四条 cut conditions
逐项验证。
\end{explanation}

\begin{exerciseblock}[Homework 7 Problem 2]
Let \(S\) be a non-empty subset of \(\R\) that has a lower bound. Prove that
\(S\) has an infimum.
\end{exerciseblock}

\begin{solution}
设 \(L\) 是 \(S\) 的 lower bound，所以对所有 \(A\in S\)，有 \(L\leq A\)。考虑
\[
  -S=\{-A:A\in S\}.
\]
Negation reverses order，因此 \(A\geq L\) 蕴含 \(-A\leq -L\)。所以 \(-S\) 是
non-empty 且 bounded above 的 subset of \(\R\)。由 Lecture Notes Exercise 48，
\(-S\) 有 supremum，记为 \(U=\sup(-S)\)。定义
\[
  I=-U.
\]
对任意 \(A\in S\)，有 \(-A\leq U\)，反转 order 得 \(I=-U\leq A\)，所以 \(I\) 是
lower bound。若 \(J\) 是 \(S\) 的任意 lower bound，则 \(J\leq A\) 对所有
\(A\in S\) 成立，故 \(-A\leq -J\) 对所有 \(A\in S\) 成立；也就是说 \(-J\) 是
\(-S\) 的 upper bound。因为 \(U\) 是 least upper bound，有 \(U\leq -J\)。再反转
order，得到 \(J\leq -U=I\)。因此 \(I\) 是 greatest lower bound，即
\(\inf S\)。
\end{solution}

\begin{explanation}
Completeness 通常只需证明 bounded-above subsets 有 supremum；infimum 可由 negative
把 lower-bound problem 转成 upper-bound problem。这个转换也解释了为什么 ordered
field 中 additive inverse 与 order 的关系很重要。同类 lower-bound 题的路线是：
先取 negatives，把 lower bounds 变成 upper bounds；取得 supremum 后再取 negative
回来，并检查 greatest lower bound 的两条定义。
\end{explanation}

\begin{exerciseblock}[Homework 7 Problem 3]
Let \(A,B,C\in\R\). Prove \((A+B)+C=A+(B+C)\).
\end{exerciseblock}

\begin{solution}
按 Dedekind addition 的定义，
\[
  (A+B)+C
  =\{(a+b)+c:a\in A,\ b\in B,\ c\in C\}.
\]
同样，
\[
  A+(B+C)
  =\{a+(b+c):a\in A,\ b\in B,\ c\in C\}.
\]
对任意 \(a,b,c\in\Q\)，rational addition associative，所以
\[
  (a+b)+c=a+(b+c).
\]
因此左边集合的每个元素都属于右边，右边集合的每个元素也属于左边。由 extensionality
of sets，
\[
  (A+B)+C=A+(B+C).
\]
\end{solution}

\begin{explanation}
Dedekind operations 的多数 algebraic laws 都 piggy-back on rational arithmetic。
这里没有 representative ambiguity，因为 cut addition 直接由所有 rational sums 组成，
而 \(\Q\) 中的 associativity 已经成立。证明 set equality 时，不要只写表达式形似；
应说明任意一个由左边生成的 rational sum 可按 \(\Q\) 的 associativity 重新括号化，
从而也由右边生成，反向同理。
\end{explanation}

\begin{exerciseblock}[Homework 7 Problem 4]
Determine whether \(A=\{x\in\Q:x^3<5\}\) is a Dedekind cut.
\end{exerciseblock}

\begin{solution}
它是 Dedekind cut。首先 \(1\in A\)，因为 \(1^3<5\)，所以 \(A\ne\emptyset\)；
\(2\notin A\)，因为 \(2^3=8>5\)，所以 \(A\ne\Q\)。

若 \(x\in A\) 且 \(y<x\)，cube function 在 \(\Q\) 上 strictly increasing，因此
\[
  y^3<x^3<5,
\]
所以 \(y\in A\)。这证明 downward closed。

最后证明 no maximum。若 \(x\in A\)，则 \(5-x^3>0\)。取 rational \(\eps>0\) 满足
\[
  \eps<1,\qquad
  \eps(3x^2+3|x|+1)<5-x^3.
\]
则
\[
  (x+\eps)^3
  =x^3+3x^2\eps+3x\eps^2+\eps^3
  \leq x^3+\eps(3x^2+3|x|+1)
  <5.
\]
因此 \(x+\eps\in A\) 且 \(x<x+\eps\)。所以 \(A\) has no maximum。
\end{solution}

\begin{explanation}
Odd powers preserve order over rationals，所以 downward closure 与 cube inequality
配合得很好。No maximum 需要一个 quantitative perturbation：选择足够小的
\(\eps\)，让 cube 的增加量小于剩余空间 \(5-x^3\)。遇到 \(x^k<c\) 的 cut 题时，
odd \(k\) 的 downward closure 通常来自 monotonicity；no maximum 则来自展开
\((x+\eps)^k\) 并把新增项压进剩余 gap。
\end{explanation}

\section{Decimal expansions and irrational numbers}

Dedekind cut 可以产生 decimal expansion。设 \(A\in\R\)。先选择 integer part，使
\[
  m_0^{\R}\leq A<(m_0+1)^{\R}.
\]
然后逐位选择 decimal truncation：对每个 \(n\geq1\)，在 denominator \(10^n\) 的
rationals 中取最大的 \(m_n/10^n\) 使它仍位于 \(A\) 的左侧。这样得到 nested
decimal intervals，长度为 \(10^{-n}\)，逐位确定 decimal digits。Terminating
expansion 和 eventually all \(9\) 的 expansion 表示同一个 cut，例如
\(1.000\ldots=0.999\ldots\)。

这里的 construction 是一个 repeated bounding process。第 \(n\) 步不是猜第
\(n\) 位 digit，而是在 mesh size \(10^{-n}\) 的 rational grid 上夹住 \(A\)：
左端点属于 cut，右端点不属于 cut。因为 cut 是 downward closed，grid 上所有更小的
点都在左边，所有足够大的点都在右边，所以可以找到最后一个仍在左边的 grid point。
相邻步骤的 intervals nested，表示每次只是在前一次的一个十分之一长度区间内继续
细分。

\begin{exerciseblock}[Lecture Notes Exercise 49]
Given a Dedekind cut, explain how to produce a decimal expansion for the
corresponding real number.
\end{exerciseblock}

\begin{solution}
设 \(A\) 是 Dedekind cut。先用 Archimedean property for rationals 找到 integer
\(m_0\) 使
\[
  m_0^{\R}\leq A<(m_0+1)^{\R}.
\]
这确定 integer part。对 \(n=1,2,3,\ldots\)，在所有 denominator 为 \(10^n\) 的
rationals 中，取 integer \(m_n\) 满足
\[
  \frac{m_n}{10^n}\in A,\qquad
  \frac{m_n+1}{10^n}\notin A.
\]
这样的 \(m_n\) 存在，因为 \(A\) 在左右都有 rationals，并且 denominator 固定后
形成 discrete grid。于是
\[
  \frac{m_n}{10^n}\leq A<\frac{m_n+1}{10^n}
\]
给出长度 \(10^{-n}\) 的 decimal interval。相邻的 \(m_n\) 彼此兼容，因此可以从
\(m_n\) 的十进制表示读出 digits。

若某一步正好落在 rational boundary，会出现两种 decimal expressions；例如某个
finite expansion 也可写成 eventually all \(9\) 的 expansion。为了得到 canonical
表达，可以规定不用 eventually all \(9\) 的形式。
\end{solution}

\begin{explanation}
Decimal expansion 是 Dedekind cut 的后续读法，而不是 definition 的基础。每一位
digit 来自一个更细的 rational grid；cut 告诉我们 grid point 应该放在左边还是右边。
遇到 decimal construction 题时，先写出 integer part 的夹逼，再写第 \(n\) 位的
grid inequality，最后说明 terminating/everything-9 ambiguity 怎样处理。
\end{explanation}

\section{Sequences and limits}

\begin{definition}[Sequence and limit]
一个 sequence in \(X\) 是 function \(\N\to X\)，通常写作 \((x_n)\)。Real sequence
\((x_n)\) has limit \(L\in\R\)，若对每个 \(\eps>0\)，存在 \(N\in\N\)，使得
\[
  n>N \Longrightarrow |x_n-L|<\eps.
\]
\end{definition}

Limit definition 的重点是 tail control。给定 tolerance \(\eps\) 后，\(N\) 只需要
控制 \(N\) 之后的所有 terms；前面有限多个 terms 可以任意大、任意震荡。写
\(\eps\)-\(N\) proof 时，先把想要的 inequality \(|x_n-L|<\eps\) 化成对 \(n\) 的
要求，再选择一个满足该要求的 \(N\)。若要证明不存在 limit，常用方法是找一个固定
\(\eps>0\)，让任何 tail 都会违反 \(|x_n-L|<\eps\)，或先证明 convergent sequence
必须 bounded 再找 contradiction。

\begin{exerciseblock}[Lecture Notes Exercise 51]
Write the definition of convergence more symbolically.
\end{exerciseblock}

\begin{solution}
Sequence \((x_n)\) converges to \(L\) 的 symbolic form 是
\[
  \forall \eps\in\R\,
  \Bigl(\eps>0\Rightarrow
  \exists N\in\N\,
  \forall n\in\N\,
  (n>N\Rightarrow |x_n-L|<\eps)\Bigr).
\]
也可以写作
\[
  \lim_{n\to\infty}x_n=L.
\]
两个写法的内容相同：每一个 positive tolerance \(\eps\) 都要求一个 tail index
\(N\)，使 tail 中所有 terms 都落在 interval \((L-\eps,L+\eps)\) 内。
\end{solution}

\begin{explanation}
符号里的 quantifier order 很重要。先给 \(\eps\)，再选择 \(N\)，最后检查所有
\(n>N\)。\(N\) 可以依赖 \(\eps\)，但不能依赖 tail 中的某个特定 \(n\)。同类题中，
不要把 ``there exists \(N\)'' 放在 ``for every \(\eps\)'' 前面；那会要求一个
single \(N\) 同时服务所有 tolerances，含义强得多。
\end{explanation}

\begin{exerciseblock}[Lecture Notes Exercise 52]
Compute the limit, if it exists, of the sequence \(x_n=0\) for all \(n\geq0\).
\end{exerciseblock}

\begin{solution}
这个 sequence 的 limit 是 \(0\)。给定任意 \(\eps>0\)，取 \(N=0\)。则对所有
\(n>N\)，
\[
  |x_n-0|=|0-0|=0<\eps.
\]
所以按 \(\eps\)-\(N\) definition，有
\[
  \lim_{n\to\infty}x_n=0.
\]
若 \(L\ne0\)，取 \(\eps=|L|/2>0\)。对所有 \(n\)，
\[
  |x_n-L|=|L|>\eps,
\]
所以 \(L\) 不能是 limit。因此 limit unique 且等于 \(0\)。
\end{solution}

\begin{explanation}
Constant zero sequence 是 definition 的 base example。证明时不需要让 \(N\) 变大，
因为所有 terms 从一开始就已经在任何以 \(0\) 为中心的 positive radius 内。若题目
还要求 uniqueness of limit，可以用 \(\eps=|L|/2\) 或更一般的
\(\eps=|L_1-L_2|/3\) 把两个不同候选 limits 分开。
\end{explanation}

\begin{exerciseblock}[Worksheet 7 Exercise 1]
Using triangle inequality, prove the listed inequalities for all \(x,y,z\in\R\).
\end{exerciseblock}

\begin{solution}
由 triangle inequality，
\[
  |x-z|=|(x-y)+(y-z)|\leq |x-y|+|y-z|.
\]
移项后得到
\[
  |x-y|+|y-z|\geq |x-z|.
\]

再证明 reverse triangle inequality 的一边。因为
\[
  x=(x-y)+y,
\]
所以
\[
  |x|\leq |x-y|+|y|.
\]
两边减去 \(|y|\)，得到
\[
  |x|-|y|\leq |x-y|.
\]
若交换 \(x,y\)，还可得到 \(|y|-|x|\leq|x-y|\)，合并即
\[
  \bigl||x|-|y|\bigr|\leq |x-y|.
\]
\end{solution}

\begin{explanation}
Triangle inequality 常用于把一个差拆成两段。第一问把 \(x-z\) 拆成
\((x-y)+(y-z)\)；第二问把 \(x\) 拆成 \((x-y)+y\)，再把 \(|y|\) 移到另一边。
做 absolute-value inequality 时，先决定要把哪个 expression 写成两项之和，再套
\(|u+v|\leq |u|+|v|\)。Reverse triangle inequality 的常用目标是把
\(\bigl||x|-|y|\bigr|\) 控制到 \(|x-y|\)。
\end{explanation}

\begin{exerciseblock}[Worksheet 7 Exercise 2]
Check whether the three listed sequences have limits, and compute the limits when
they exist.
\end{exerciseblock}

\begin{solution}
\begin{enumerate}[label=(\alph*)]
  \item
  \[
    x_n=\frac{5n^2+n+7}{n^3+3n^2+3n+1}.
  \]
  对 \(n\geq1\)，denominator positive，并且
  \[
    0\leq x_n
    \leq \frac{13n^2}{n^3}
    =\frac{13}{n}.
  \]
  给定 \(\eps>0\)，取 \(N>13/\eps\)，则 \(n>N\) 时 \(|x_n-0|<\eps\)。所以
  \(x_n\to0\)。

  \item
  \[
    x_n=\frac{n^2-(-1)^n n-13}{n^2+n+1}.
  \]
  计算
  \[
    x_n-1
    =\frac{-(-1)^n n-13-n-1}{n^2+n+1}.
  \]
  因此
  \[
    |x_n-1|
    \leq\frac{|(-1)^n|n+n+14}{n^2}
    =\frac{2n+14}{n^2}
    =\frac2n+\frac{14}{n^2}.
  \]
  给定 \(\eps>0\)，取 \(N\) 使 \(2/N<\eps/2\) 且 \(14/N^2<\eps/2\)。当
  \(n>N\) 时，\(|x_n-1|<\eps\)。所以 \(x_n\to1\)。

  \item \(x_n=2^n\) 没有 finite real limit。若 \(x_n\to L\)，则 convergent
  sequence bounded。取 \(M>|L|+1\) 作为 tail bound，再把有限多个 early terms 合并，
  会得到所有 \(2^n\) 的共同 upper bound。但对任意 real \(M\)，Archimedean
  property 给出 \(n\) 使 \(2^n>M\)，矛盾。因此 \(2^n\) diverges。
\end{enumerate}
\end{solution}

\begin{explanation}
Rational functions 的 limit 由最高次项控制，但 proof 仍要写成 inequality。第二问
的 oscillating term \((-1)^n n\) 只是一阶误差；除以 \(n^2\) 后趋于 \(0\)。第三问
用 boundedness 排除 convergence。处理类似 rational sequence 时，可以先除以最高
次幂，或直接把 numerator 用 \(Cn^k\) 估计；目标是制造 \(C/n^r\) 这样的 upper
bound，再用 \(N>C/\eps\) 收尾。
\end{explanation}

\begin{exerciseblock}[Worksheet 7 Exercise 3]
Let \((x_n)\) be a sequence of real numbers with limit \(L\). Prove that it is
bounded.
\end{exerciseblock}

\begin{solution}
因为 \(x_n\to L\)，对 \(\eps=1\)，存在 \(N\in\N\) 使 \(n>N\) 时
\[
  |x_n-L|<1.
\]
这意味着 tail 中所有 terms 都满足
\[
  L-1<x_n<L+1.
\]
前面只有有限多个 terms：\(x_0,x_1,\ldots,x_N\)。令
\[
  m=\min\{x_0,x_1,\ldots,x_N,L-1\},
  \qquad
  M=\max\{x_0,x_1,\ldots,x_N,L+1\}.
\]
有限集合在 total order 中有 maximum 和 minimum。若 \(n\leq N\)，则
\(m\leq x_n\leq M\) 由定义成立；若 \(n>N\)，则 \(L-1<x_n<L+1\)，也有
\[
  m\leq x_n\leq M.
\]
所以存在 \(m,M\in\R\) 使所有 \(n\) 都满足 \(m\leq x_n\leq M\)。
\end{solution}

\begin{explanation}
Convergence 只直接控制 tail；boundedness 需要同时处理有限多个 early terms。有限
多个 early terms 可以取 maximum 和 minimum，因此不会破坏整体 boundedness。同类题
的结构固定：用 \(\eps=1\) 控制 tail，再把 finite initial segment 和 tail bound
放进同一个 maximum/minimum 中。
\end{explanation}

\begin{exerciseblock}[Worksheet 7 Exercise 4]
If \(x_n\to L\), prove \(|x_n|\to |L|\).
\end{exerciseblock}

\begin{solution}
由 Worksheet 7 Exercise 1 的 reverse triangle inequality，
\[
  \bigl||x_n|-|L|\bigr|\leq |x_n-L|.
\]
给定 \(\eps>0\)，因为 \(x_n\to L\)，存在 \(N\) 使 \(n>N\) 时
\[
  |x_n-L|<\eps.
\]
于是同样的 \(N\) 满足
\[
  \bigl||x_n|-|L|\bigr|<\eps.
\]
所以按 limit definition，
\[
  \lim_{n\to\infty}|x_n|=|L|.
\]
\end{solution}

\begin{explanation}
Absolute value function 的 continuity 在这里由 reverse triangle inequality 表达。
它把新 sequence \(|x_n|\) 与原 sequence \(x_n\) 的距离直接比较起来。证明由已知
limit 推出新 limit 时，最理想的形式是
\[
  |\,\text{new error}\,|\leq C|\,\text{old error}\,|
\]
或更强的 \(|\,\text{new error}\,|\leq |\,\text{old error}\,|\)，然后沿用原来的
\(N\)。
\end{explanation}

\begin{exerciseblock}[Worksheet 7 Exercise 5]
Calculate \(\lim_{n\to\infty}(\sqrt{n+1}-\sqrt n)\).
\end{exerciseblock}

\begin{solution}
对每个 \(n\geq1\)，有
\[
  \sqrt{n+1}-\sqrt n
  =\frac{(\sqrt{n+1}-\sqrt n)(\sqrt{n+1}+\sqrt n)}
         {\sqrt{n+1}+\sqrt n}
  =\frac1{\sqrt{n+1}+\sqrt n}.
\]
因为 denominator 大于 \(\sqrt n\)，所以
\[
  0<\sqrt{n+1}-\sqrt n<\frac1{\sqrt n}.
\]
给定 \(\eps>0\)，取 \(N>1/\eps^2\)。若 \(n>N\)，则 \(\sqrt n>1/\eps\)，从而
\[
  0<\sqrt{n+1}-\sqrt n<\eps.
\]
因此 limit 为 \(0\)。
\end{solution}

\begin{explanation}
两个 square roots 的差不适合直接估计；乘以 conjugate 后变成 reciprocal。分母随
\(n\) 增大而趋向 infinity，所以整个 expression 趋向 \(0\)。含 radical 的 limit
题常先尝试 rationalization；若目标是证明趋于 \(0\)，把表达式压到
\(1/\sqrt n\) 或 \(1/n\) 这样的简单 upper bound 上即可。
\end{explanation}

\section{Cauchy sequences \optional}

Lecture Notes 4.13 给出 real numbers 的另一种 construction：从 rational Cauchy
sequences 出发，把彼此最终无限接近的 sequences 视为同一个 real number。

这个 construction 的顺序是：先收集所有 rational Cauchy sequences；再定义
equivalence relation
\[
  (x_n)\sim(y_n)
  \Longleftrightarrow
  \forall\eps>0\ \exists N\ \forall n>N,\ |x_n-y_n|<\eps;
\]
最后把 equivalence classes 当作 real numbers。和 Dedekind cuts 相比，这里不再把
real number 看成 rational line 的左半边，而是看成一类会逐渐稳定到同一位置的
rational approximations。

\begin{definition}[Cauchy sequence]
Rational sequence \((x_n)\) 称为 Cauchy，若
\[
  \forall\eps>0\ \exists N\in\N\ \forall m,n>N,\quad |x_n-x_m|<\eps.
\]
这个 condition 不提未知的 limit，只要求 tail 中任意两个 terms 最终彼此接近。
\end{definition}

Cauchy condition 与 convergence definition 的量词很相似，但最后比较的是
\(|x_n-x_m|\)，不是 \(|x_n-L|\)。在 complete spaces 中，Cauchy 等价于
convergent；在 \(\Q\) 中它更有用，因为像 decimal approximations to \(\sqrt2\)
这样的 rational sequence 可以 Cauchy，却没有 rational limit。Cauchy construction
正是把这些 missing limits 补成新的 elements。

\begin{exerciseblock}[Lecture Notes Exercise 53]
Find examples of Cauchy sequences and sequences which fail to be Cauchy.
\end{exerciseblock}

\begin{solution}
Constant sequence \(x_n=q\) 是 Cauchy。给定 \(\eps>0\)，任取 \(N\)，则
\[
  |x_n-x_m|=|q-q|=0<\eps
\]
对所有 \(m,n>N\) 成立。

Sequence \(x_n=1/n\) 也是 Cauchy。给定 \(\eps>0\)，取 \(N>2/\eps\)。若
\(m,n>N\)，则
\[
  \left|\frac1n-\frac1m\right|
  \leq \frac1n+\frac1m
  <\frac1N+\frac1N
  <\eps.
\]

Sequence \(x_n=n\) 不是 Cauchy。取 \(\eps=1\)。无论 \(N\) 多大，选择
\(n=N+2\)、\(m=N+1\)，都有
\[
  |x_n-x_m|=1,
\]
不小于 \(\eps\)。Sequence \(x_n=(-1)^n\) 也不是 Cauchy；取 \(\eps=1\)，tail 中
总有一偶一奇两个 indices，使对应 terms 相差 \(2\)。
\end{solution}

\begin{explanation}
Cauchy condition 描述的是 tail 的 internal closeness。Converging-to-zero 的
\(1/n\) 满足它；持续向 infinity 走的 \(n\) 和来回震荡的 \((-1)^n\) 都不满足。
构造 non-Cauchy example 时，不必知道它有没有 limit；只要找到一个 fixed
\(\eps\)，并证明每个 tail 里仍能找到两项距离至少 \(\eps\)，就已经否定 Cauchy。
\end{explanation}

\begin{exerciseblock}[Lecture Notes Exercise 54]
Explain how rationals appear in the Cauchy-sequence construction of \(\R\).
\end{exerciseblock}

\begin{solution}
每个 rational \(q\in\Q\) 对应 constant Cauchy sequence
\[
  (q,q,q,\ldots).
\]
定义 embedding
\[
  j:\Q\to\R,\qquad j(q)=[(q,q,q,\ldots)].
\]
它 well-defined，因为 constant sequence 是 Cauchy。它 injective：若
\[
  [(p,p,p,\ldots)]=[(q,q,q,\ldots)],
\]
则按 equivalence definition，对每个 \(\eps>0\)，eventually 有
\[
  |p-q|<\eps.
\]
若 \(p\ne q\)，取 \(\eps=|p-q|/2>0\)，就得到 \(|p-q|<|p-q|/2\)，矛盾。因此
\(p=q\)。
\end{solution}

\begin{explanation}
Cauchy construction 中的 real number 是 equivalence class。Rationals 进入这个
quotient set 的方式是 constant representatives；不同 rationals 的 constant
sequences 不会被 equivalence relation 合并。证明 embedding injective 的关键是
拿 \(\eps=|p-q|/2\)；这会把两个不同 constant sequences 永久分开。
\end{explanation}

\begin{exerciseblock}[Lecture Notes Exercise 55]
Find a few representatives of \(1/2\in\Q\) in the Cauchy-sequence construction.
\end{exerciseblock}

\begin{solution}
以下 rational Cauchy sequences 都代表 \(1/2\)：
\[
  x_n=\frac12,\qquad
  y_n=\frac12+\frac1n,\qquad
  z_n=\frac12+\frac{(-1)^n}{n}.
\]
它们都是 Cauchy，因为 constant sequence 是 Cauchy，\(1/n\) 是 Cauchy，且 Cauchy
sequences 的 sum 与 sign change 仍是 Cauchy。

还要说明它们与 constant representative equivalent。对 \(y_n\)，给定
\(\eps>0\)，取 \(N>1/\eps\)。若 \(n>N\)，则
\[
  \left|y_n-\frac12\right|=\frac1n<\eps.
\]
对 \(z_n\)，同样取 \(N>1/\eps\)，则
\[
  \left|z_n-\frac12\right|=\frac1n<\eps.
\]
所以这些 sequences 都属于 \(1/2\) 的 equivalence class。
\end{solution}

\begin{explanation}
同一个 real number 可以有很多 representatives。只要 sequence 的 error 相对于
\(1/2\) 趋向 \(0\)，它就在 constant \(1/2\) sequence 的 equivalence class 中。
寻找 representatives 的方法是从 constant sequence 出发，加上任何趋于 \(0\) 的
rational sequence；例如 \(1/n\)、\((-1)^n/n\)、或 \(1/n^2\) 都只改变代表而不改变
equivalence class。
\end{explanation}

\begin{exerciseblock}[Lecture Notes Exercise 56]
Show that every nonzero Cauchy-sequence real number has a multiplicative inverse.
\end{exerciseblock}

\begin{solution}
设 \(x=[(x_n)]\ne[(0,0,0,\ldots)]\)。Nonzero 的意思是 \((x_n)\) 不 equivalent to
zero sequence。因此存在 \(\eta>0\)，使得对每个 \(N\)，都能找到 \(n>N\) 满足
\(|x_n|\geq\eta\)；否则 \(|x_n|\) 对每个 positive tolerance 最终都小于该
tolerance，\((x_n)\) 就会 equivalent to zero。

现在证明存在 \(N_0\) 使所有 \(n>N_0\) 都满足
\[
  |x_n|\geq\eta/2.
\]
因为 \((x_n)\) 是 Cauchy，取 \(N_1\) 使 \(m,n>N_1\) 时
\[
  |x_m-x_n|<\eta/2.
\]
由上一段，存在 \(k>N_1\) 满足 \(|x_k|\geq\eta\)。对任意 \(n>N_1\)，triangle
inequality 给出
\[
  |x_n|\geq |x_k|-|x_n-x_k|>\eta-\eta/2=\eta/2.
\]
所以可取 \(N_0=N_1\)。

定义 rational sequence
\[
  y_n=
  \begin{cases}
    0,& n\leq N_0,\\
    1/x_n,& n>N_0.
  \end{cases}
\]
若 \(m,n>N_0\)，则
\[
  |y_n-y_m|
  =\left|\frac1{x_n}-\frac1{x_m}\right|
  =\frac{|x_m-x_n|}{|x_nx_m|}
  \leq \frac{4}{\eta^2}|x_m-x_n|.
\]
因为 \((x_n)\) is Cauchy，右边可任意小，所以 \((y_n)\) is Cauchy。

最后，对 \(n>N_0\)，有 \(x_ny_n=1\)。因此 product sequence \((x_ny_n)\) 与
constant sequence \((1,1,1,\ldots)\) 只在有限多个位置可能不同，二者 equivalent。
于是
\[
  [(x_n)]\cdot[(y_n)]=[(1,1,1,\ldots)].
\]
所以 \(x\) 有 multiplicative inverse。
\end{solution}

\begin{explanation}
核心难点是避免除以接近 \(0\) 的 terms。Nonzero equivalence class 保证某个 tail
远离 \(0\)；有了这个 lower bound，reciprocal sequence 才会保持 Cauchy。同类
``take reciprocal'' proof 都要先建立 uniform lower bound \(|x_n|\geq c>0\)，再用
\[
  \left|\frac1{x_n}-\frac1{x_m}\right|
  =\frac{|x_m-x_n|}{|x_nx_m|}
\]
把 reciprocal error 控制回原 sequence 的 Cauchy error。
\end{explanation}

\section{Homework 8: ordered-field consequences}

这一节把前面构造出的 \(\R\) 当作 complete ordered field 使用。证明风格从
``某个 set 是不是 Dedekind cut'' 转向 ``ordered-field axiom 加 completeness 能推出
哪些分析性质''。读这些题时要分清三种工具：trichotomy 来自 total order，inverse
来自 field operation 的构造，Archimedean property 和 square-root theorem 则需要
least-upper-bound argument。

\begin{exerciseblock}[Homework 8 Problem 1]
Trichotomy. Let \(A\in\R\). Prove exactly one of \(A>0^{\R}\), \(A=0^{\R}\),
or \(A<0^{\R}\) holds.
\end{exerciseblock}

\begin{solution}
Order on Dedekind cuts is total，因此对 \(A\) 与 \(0^{\R}\)，有
\[
  A\leq0^{\R}\quad\text{or}\quad0^{\R}\leq A.
\]
若两者同时成立，则 \(A\subseteq0^{\R}\) 且 \(0^{\R}\subseteq A\)，所以
\(A=0^{\R}\)。

若 \(0^{\R}\leq A\) 且 \(A\ne0^{\R}\)，则按 strict order definition，
\[
  A>0^{\R}.
\]
若 \(A\leq0^{\R}\) 且 \(A\ne0^{\R}\)，则
\[
  A<0^{\R}.
\]
三种情况至少有一种成立。

它们不能同时成立：\(A>0^{\R}\) 与 \(A<0^{\R}\) 会给出
\[
  0^{\R}<A<0^{\R},
\]
违反 antisymmetry/transitivity；任一 strict inequality 也排除 equality。因此
exactly one holds。
\end{solution}

\begin{explanation}
Trichotomy 是 total order 的具体表现。先用 totality 得到可比较性，再用
antisymmetry 把 double inclusion 识别成 equality，剩下的 proper inclusion 就是
strict inequality。证明 ``exactly one'' 要包含两部分：至少一个成立，以及不能有两个
同时成立。
\end{explanation}

\begin{exerciseblock}[Homework 8 Problem 2]
Multiplicative inverse. Let \(A\in\R\), \(A\ne0^{\R}\). Prove there exists a
unique \(B\in\R\) such that \(A\cdot B=1^{\R}\).
\end{exerciseblock}

\begin{solution}
先处理 \(A>0^{\R}\)。定义
\[
  B=\{q\in\Q:q<0\}\cup
    \{q\in\Q:q>0\text{ and }q<1/c
      \text{ for some }c\in\Q\setminus A,\ c>0\}.
\]
这个 \(B\) non-empty，因为包含所有 negative rationals；proper，因为若
\(a\in A\) 且 \(a>0\)，则任何 \(q\geq1/a\) 都不能属于 \(B\)。Downward closed
由 inequality \(u<q<1/c\) 继承。No maximum 分两类：negative element 可向右移动
仍保持 negative，positive element \(q<1/c\) 可在 \(q\) 与 \(1/c\) 之间取 rational
\(q'\)，仍满足 \(q'<1/c\)。所以 \(B\) 是 Dedekind cut。

按照 positive cuts 的 multiplication definition，右半边的 elements \(c\notin A\)
和 \(d\notin B\) 都是 positive boundary-side rationals。由 \(d\notin B\) 可得
\[
  d\geq 1/c
\]
对每个 positive \(c\notin A\) 成立，所以 \(cd\geq1\)。这说明所有 \(x<1\) 都在
\(A\cdot B\)，即 \(1^{\R}\subseteq A\cdot B\)。

反过来，先处理 \(x>1\)。选 positive \(a\in A\)。因为 \(x>1\)，有 \(a<xa\)。
右半边 \(\Q\setminus A\) 非空，且 \(A\) 没有 maximum；用 rational density 可在
boundary 右侧选 \(c\notin A\) 使
\[
  c/x\in A.
\]
令 \(d=x/c\)。若 \(e\in\Q\setminus A\) 且 \(e>0\)，则不能有 \(e<c/x\)，否则
downward closed 会由 \(c/x\in A\) 推出 \(e\in A\)。所以 \(e\geq c/x\)，从而
\[
  d=\frac{x}{c}\geq \frac1e.
\]
这对每个 positive \(e\notin A\) 成立，故 \(d\notin B\)。同时 \(c\notin A\) 且
\[
  cd=x.
\]
按照 product cut 的 upper-side characterization，存在 \(c\notin A,d\notin B\)
且 \(cd\leq x\)，所以 \(x\notin A\cdot B\)。于是所有 \(x>1\) 都不在 product cut。
若 \(1\in A\cdot B\)，由于 \(A\cdot B\) has no maximum，会存在 \(y\in A\cdot B\)
使 \(y>1\)，这与刚证明的结论矛盾。因此 \(1\notin A\cdot B\)。综上，
\[
  A\cdot B=1^{\R}.
\]

若 \(A<0^{\R}\)，则 \(-A>0^{\R}\)。按上半部分构造 \(C\) 使
\[
  (-A)C=1^{\R}.
\]
令 \(B=-C\)。由 sign rule for multiplication，
\[
  A B=A(-C)=(-A)C=1^{\R}.
\]

Uniqueness：若 \(AB=1^{\R}\) 且 \(AC=1^{\R}\)，则
\[
  B=B\cdot1^{\R}=B(AC)=(BA)C=(AB)C=1^{\R}C=C.
\]
因此 inverse unique。
\end{solution}

\begin{explanation}
Positive case 的 construction 把 \(A\) 的右半边 \(c\) 反转成 \(1/c\) 的左半边。
负数情况不用重新构造，只需先取 additive inverse 变成 positive cut，再把 inverse
的 sign 反回来。做 inverse theorem 时，不要一开始就写 \(1/A\)，因为 \(A\) 还不是
ordinary number；必须用 cut 的右半边描述哪些 rational reciprocals 应该落在左边。
Uniqueness 则是 field-style cancellation：把两个候选 inverse 先乘上同一个 \(A\)，
再用 associativity 和 identity 化简。
\end{explanation}

\begin{exerciseblock}[Homework 8 Problem 3]
Archimedean property. Let \(A,B\in\R\) be positive cuts. Prove that there exists
an integer \(N\) such that \(N^{\R}\cdot A>B\).
\end{exerciseblock}

\begin{solution}
反设对所有 \(N\in\N\)，都有
\[
  N^{\R}\cdot A\leq B.
\]
令
\[
  T=\{N^{\R}\cdot A:N\in\N\}.
\]
则 \(T\) non-empty 且 bounded above by \(B\)。由 completeness，存在
\[
  C=\sup T.
\]
因为 \(A>0^{\R}\)，有 \(C-A<C\)。若 \(C-A\) 仍是 \(T\) 的 upper bound，则由于
\(C\) 是 least upper bound，应有 \(C\leq C-A\)，这与 \(A>0^{\R}\) 矛盾。因此
\(C-A\) 不是 upper bound。于是存在 \(N\in\N\) 使
\[
  N^{\R}\cdot A>C-A.
\]
两边加 \(A\)，得
\[
  (N+1)^{\R}\cdot A>C.
\]
但 \((N+1)^{\R}\cdot A\in T\)，而 \(C\) 是 \(T\) 的 upper bound，矛盾。因此原
假设不成立，存在 \(N\) 使 \(N^{\R}\cdot A>B\)。
\end{solution}

\begin{explanation}
Archimedean property 的证明使用 completeness。如果所有 multiples 都被 \(B\) 压住，
它们就有 supremum \(C\)。但从 \(C\) 往左减去一个 positive \(A\) 后不可能仍是
upper bound，于是某个 multiple 超过 \(C-A\)，再加一次 \(A\) 就超过 \(C\)。
同类 proof by contradiction 的特征是：假设某个 increasing family 永远被 bounded
above，取 supremum，然后利用 ``supremum 之前一点不是 upper bound'' 找到超过它的
下一步。
\end{explanation}

\begin{exerciseblock}[Homework 8 Problem 4]
Square root. Let \(A\in\R\) be non-negative. Prove there exists a unique
non-negative \(B\in\R\) such that \(B\cdot B=A\).
\end{exerciseblock}

\begin{solution}
若 \(A=0^{\R}\)，取 \(B=0^{\R}\) 即可。下面设 \(A>0^{\R}\)。定义
\[
  S=\{X\in\R: X\geq0^{\R}\text{ and }X^2\leq A\}.
\]
这个 set non-empty，因为 \(0^{\R}\in S\)。它 bounded above：由 Archimedean
property，可取 integer \(N\) 使 \(N^{\R}>A+1^{\R}\)。若 \(X\geq N^{\R}\)，则
\(X^2\geq N^{\R}X\geq N^{\R}>A\)，所以这样的 \(X\) 不在 \(S\)。因此 \(N^{\R}\)
是 upper bound。由 completeness，令
\[
  B=\sup S.
\]
则 \(B\geq0^{\R}\)。

证明 \(B^2=A\)。若 \(B^2<A\)，令差为 \(D=A-B^2>0\)。取 small positive rational
\(h^{\R}\) 使
\[
  2Bh+h^2<D.
\]
则
\[
  (B+h)^2=B^2+2Bh+h^2<A,
\]
所以 \(B+h\in S\)，但 \(B+h>B\)，与 \(B\) 是 upper bound 矛盾。

若 \(B^2>A\)，令 \(D=B^2-A>0\)。取 small positive rational \(h^{\R}\) 满足
\[
  0<h<B,\qquad 2Bh-h^2<D.
\]
则
\[
  (B-h)^2=B^2-2Bh+h^2>A.
\]
若 \(X\in S\)，则 \(X^2\leq A<(B-h)^2\)。因为 \(X,B-h\) non-negative，order 与
multiplication 相容给 \(X\leq B-h\)。所以 \(B-h\) 是 \(S\) 的 upper bound，且
\(B-h<B\)，与 \(B=\sup S\) 的 leastness 矛盾。因此只能有 \(B^2=A\)。

Uniqueness：若 \(B,C\geq0^{\R}\) 且 \(B^2=C^2=A\)。若 \(B<C\)，则
\[
  0^{\R}<C-B,\qquad 0^{\R}\leq B+C,
\]
并且 \(B+C>0^{\R}\)。于是
\[
  0^{\R}<(C-B)(B+C)=C^2-B^2=0^{\R},
\]
矛盾。若 \(C<B\)，交换 \(B\) 与 \(C\) 的位置，同一段计算给出
\[
  0^{\R}<(B-C)(B+C)=B^2-C^2=0^{\R},
\]
也矛盾。所以既不能 \(B<C\)，也不能 \(C<B\)。由 trichotomy，只能 \(B=C\)。
\end{solution}

\begin{explanation}
Square root existence 是 supremum method 的典型应用：把所有平方不超过 \(A\) 的
non-negative numbers 收集起来，取它们的 supremum。若 supremum 的平方太小，可向右
移动；若平方太大，可向左移动。两种移动都由 small rational perturbation 控制。解决
类似 existence theorem 时，先定义候选集合 \(S\)，再证明 \(S\) non-empty and
bounded above，接着令 \(B=\sup S\)，最后用 ``too small'' 和 ``too large'' 两个
contradictions 把 \(B\) 钉到目标 equation 上。
\end{explanation}

\begin{reviewbox}{Chapter checklist}
本章需要掌握：total order, ordered field, upper/lower bound, supremum/infimum,
completeness, \(\Q\) not complete, Dedekind cut 的 pair and one-set versions,
rational cut embedding, cut order, addition, negative, multiplication, completeness
by union, decimal expansion, irrational cut for \(\sqrt2\), sequence limit,
Cauchy sequence quotient construction, and ordered-field consequences such as
inverse, Archimedean property, and square roots。
\end{reviewbox}
